\documentclass[11pt]{article}
\usepackage[margin=1in]{geometry}
\usepackage[T1]{fontenc}
\usepackage{lmodern,microtype,hyperref,enumitem,amsmath,amssymb,xcolor,booktabs}
\usepackage[numbers,sort&compress]{natbib}
\usepackage{url}
\hypersetup{colorlinks=true,urlcolor=blue,citecolor=blue,linkcolor=blue}

\title{Q12 --- Comprehension check}
\author{DLS 2026 --- AI Internals Module}
\date{}

\begin{document}
\maketitle

\section*{Student question}

\emph{``Why do models give confident answers even when they are wrong?''}

This comprehension check asks you to connect the surface phenomenon of
``confident wrongness'' to concrete facts about how language models are
trained, calibrated, and decoded.

\paragraph{Question.}
A large language model answers a factual question fluently and
authoritatively, yet the stated fact is false (a ``hallucination'').
Explain \emph{why} such models tend to express high confidence even when
wrong. In your answer, address \textbf{all four} of the following
mechanisms, and state briefly how each contributes to the mismatch
between expressed confidence and actual correctness:
\begin{enumerate}[label=(\alph*)]
  \item the training objective (next-token likelihood) and what it does
        \emph{not} reward;
  \item \emph{miscalibration} introduced or reshaped by post-training
        (RLHF / instruction tuning);
  \item the effect of \emph{decoding} and token-level probabilities on
        the fluency of a wrong answer;
  \item why the model has no reliable internal signal that cleanly
        separates ``I know'' from ``I am guessing.''
\end{enumerate}

\paragraph{Answer.}
Confident wrongness is not a single bug but the joint product of the
objective, the alignment procedure, and the decoding regime.

\textbf{(a) The objective rewards plausibility, not truth.}
Autoregressive models are trained to maximize the likelihood of the next
token given the context \citep{radford2019language,brown2020language}.
This objective rewards producing the \emph{statistically most plausible
continuation} of the preceding text, including the fluent, assertive
style of confident prose that dominates the training corpus. Nothing in
the loss distinguishes a true statement from a false-but-plausible one;
a well-formed wrong answer and a well-formed right answer can incur
nearly identical per-token loss. Hallucination is therefore partly
intrinsic: \citet{kalai2024calibrated} and related analyses argue that
models optimized this way will inevitably generate some fabricated but
statistically typical content, and that standard training and evaluation
\emph{reward guessing over abstaining} \citep{kalai2025hallucinate}.

\textbf{(b) Post-training reshapes (and often degrades) calibration.}
A pretrained base model can be surprisingly well-calibrated on
multiple-choice tasks: its probabilities roughly track the likelihood of
being correct \citep{kadavath2022know}. However, RLHF and instruction
tuning optimize for human-preferred responses, and human raters tend to
prefer answers that are decisive, helpful-sounding, and confident.
OpenAI's GPT-4 report explicitly notes that the base model was
well-calibrated but that calibration \emph{decreased} after RLHF
post-training \citep{openai2023gpt4}. The model thus learns a stylistic
association between confidence and reward, decoupling expressed certainty
from epistemic warrant. Preference optimization can also induce
\emph{sycophancy}---agreeing with or flattering the user---which further
rewards assured-sounding but unjustified claims
\citep{perez2023discovering}.

\textbf{(c) Decoding produces fluent surface form regardless of truth.}
At generation time the model samples or greedily selects high-probability
tokens. Because fluency and factuality are governed by different things,
a false claim can still be composed of individually high-probability
tokens, yielding text that \emph{reads} as confident. The sequence-level
probability is a measure of linguistic typicality, not veracity, so a
hallucinated answer can be assigned high likelihood. Methods such as
semantic-entropy estimation show that meaningful uncertainty signal
exists but requires sampling multiple generations and measuring
\emph{meaning}-level disagreement rather than reading off a single
decoded string \citep{kuhn2023semantic,farquhar2024detecting}.

\textbf{(d) No clean internal ``I know / I'm guessing'' signal is
surfaced by default.}
Probing work indicates models do carry \emph{some} information about
their own reliability---they can be trained to predict whether they know
an answer \citep{kadavath2022know}---but this signal is noisy, is not
faithfully exposed in the free-text output, and is weakened by the
confidence-rewarding post-training above. Verbalized confidence
estimates are often poorly calibrated and must be elicited and corrected
explicitly \citep{lin2022teaching}. Consequently, in default use the
model has no mechanism that converts internal uncertainty into an honest
hedge; it emits the same assertive register whether the retrieved
association is strong or spurious.

\textbf{Synthesis.} The model expresses confidence because confidence is
cheap, stylistically rewarded, and statistically typical, while truth is
never directly optimized and uncertainty is never reliably surfaced.
Mitigations target exactly these points: calibration-aware training and
evaluation that reward abstention \citep{kalai2025hallucinate},
retrieval grounding, sampling-based uncertainty measures
\citep{farquhar2024detecting}, and explicit confidence elicitation
\citep{lin2022teaching}.

\bibliographystyle{plainnat}
\begin{thebibliography}{10}

\bibitem[Brown et al.(2020)]{brown2020language}
T.~B. Brown, B.~Mann, N.~Ryder, et~al.
\newblock Language models are few-shot learners.
\newblock In \emph{Advances in Neural Information Processing Systems (NeurIPS)}, 2020.

\bibitem[Farquhar et al.(2024)]{farquhar2024detecting}
S.~Farquhar, J.~Kossen, L.~Kuhn, and Y.~Gal.
\newblock Detecting hallucinations in large language models using semantic entropy.
\newblock \emph{Nature}, 630:625--630, 2024.

\bibitem[Kadavath et al.(2022)]{kadavath2022know}
S.~Kadavath, T.~Conerly, A.~Askell, et~al.
\newblock Language models (mostly) know what they know.
\newblock \emph{arXiv preprint arXiv:2207.05221}, 2022.

\bibitem[Kalai and Vempala(2024)]{kalai2024calibrated}
A.~T. Kalai and S.~S. Vempala.
\newblock Calibrated language models must hallucinate.
\newblock In \emph{Proceedings of the 56th Annual ACM Symposium on Theory of Computing (STOC)}, 2024.

\bibitem[Kalai et al.(2025)]{kalai2025hallucinate}
A.~T. Kalai, O.~Nachum, S.~S. Vempala, and E.~Zhang.
\newblock Why language models hallucinate.
\newblock \emph{arXiv preprint arXiv:2509.04664}, 2025.

\bibitem[Kuhn et al.(2023)]{kuhn2023semantic}
L.~Kuhn, Y.~Gal, and S.~Farquhar.
\newblock Semantic uncertainty: Linguistic invariances for uncertainty estimation in natural language generation.
\newblock In \emph{International Conference on Learning Representations (ICLR)}, 2023.

\bibitem[Lin et al.(2022)]{lin2022teaching}
S.~Lin, J.~Hilton, and O.~Evans.
\newblock Teaching models to express their uncertainty in words.
\newblock \emph{Transactions on Machine Learning Research (TMLR)}, 2022.

\bibitem[OpenAI(2023)]{openai2023gpt4}
OpenAI.
\newblock {GPT-4} technical report.
\newblock \emph{arXiv preprint arXiv:2303.08774}, 2023.

\bibitem[Perez et al.(2023)]{perez2023discovering}
E.~Perez, S.~Ringer, K.~Luko\v{s}i\=ut\.e, et~al.
\newblock Discovering language model behaviors with model-written evaluations.
\newblock In \emph{Findings of the Association for Computational Linguistics (ACL)}, 2023.

\bibitem[Radford et al.(2019)]{radford2019language}
A.~Radford, J.~Wu, R.~Child, D.~Luan, D.~Amodei, and I.~Sutskever.
\newblock Language models are unsupervised multitask learners.
\newblock Technical report, OpenAI, 2019.

\end{thebibliography}

\end{document}